\textbf{COMP6008 \textbar{} Group 7 \textbar{} Working research dossier for independent report preparation}

\textbf{NOT A SUBMISSION DOCUMENT --- contributors and report text must be independently prepared under the unit's no-generative-AI policy.}

Reviewed against July 29, 2026 guidelines and repository master at b9f6b7438afafcb25455f707ecda7b004d6fc457.

\textbf{ABSTRACT---Evidence synopsis.} A single centrally controlled fleet of three drones performs recurring inspections of simulated wildfire hazards while managing shared charging, battery depletion, evolving fire, and inspection deadlines. DQN, REINFORCE, and synchronous A2C are implemented and evaluated in a fully observable, finite-horizon Markov decision process. A frozen 40$\times$40 comparison comprises nine independent trained policies and 260 test missions. Greedy achieves 0.22077 mean weighted freshness, versus DQN 0.08438, REINFORCE 0.06901, and A2C 0.08435. Masked random yields 0.08379. Zero losses for masked learners therefore evidence strong rule-based feasibility control, not a demonstrated learned surveillance advantage. The notes below identify verified methods, limitations, and literature-grounded follow-ups; they are for independent student synthesis, not direct submission.

\textbf{INDEX TERMS---}reinforcement learning; wildfire monitoring; Markov decision processes; drone coordination; DQN; policy gradient; actor--critic; policy evaluation.
